\clearpage
\renewcommand{\BookRunningTitle}{Geometrische Seite der ternären Formen}
\SourceLeaf{315}
\pdfbookmark[0]{Geometrische Seite der ternären Formen}{ternary-geometry}
\begin{center}\Large\bfseries GEOMETRISCHE SEITE DER TERNÄREN FORMEN.\end{center}
\BookRule{30mm}

Ein Punkt im Raum $(0)$ werde als Anfangspunkt gewählt. Der Übergang von da zu drei anderen Punkten $P,P',P''$, die mit jenem nicht in einer Ebene liegen, sei jeweils $t,t',t''$; wobei, wenn keine Verwechslung möglich ist, die Punkte $P,P',P''$ selbst durch $(t)$. $(t')$, $(t'')$ bezeichnet werden können.

Ferner sei allgemein $(t,t')$ das Produkt der Längen der beiden Linien $t,t'$ mit dem Kosinus ihres Neigungswinkels usw.

Man hat allgemein $(\alpha t+\alpha't'+\alpha''t''+\ldots\quad\beta u+\beta'u'+\beta''u''+\mathinner{.\,.})$ wenn man die Multiplikation
\[
(\alpha t+\alpha't'+\alpha''t''+\mathinner{.\,.})\times(\beta u+\beta'u'+\beta''u''+\mathinner{.\,.})
\]
ausführt und statt $tu,\allowbreak tu',\allowbreak tu'',\allowbreak t'u',\allowbreak t'u''\mathinner{.\,.}$ usw. $(t,u),\allowbreak(t,u'),\allowbreak(t,u''),\allowbreak(t',u'),\allowbreak(t',u'')$ usw. schreibt.

Jeder Punkt im Raum kann durch ein Trinom
\[
(xt+x't'+x''t'')
\]
dargestellt werden.

Für alle Punkte, die in einer bestimmten Ebene liegen, wird dann eine Gleichung
\[
\lambda x+\lambda'x'+\lambda''x''=L
\]
\SourceLeaf{316}
gelten, wo $\lambda,\lambda',\lambda'',L$ feste Zahlen bezeichnen. Für eine Ebene durch die drei Punkte $\mu t,\mu't',\mu''t''$ ist
\[
\lambda\mu=\lambda'\mu'=\lambda''\mu''=L
\]

Schreibt man
\[
\begin{gathered}
(t,t)=a,\quad(t',t')=a',\quad(t'',t'')=a'',\\
(t',t'')=b,\quad(t,t'')=b',\quad(t,t')=b''
\end{gathered}
\]
und
\[
\begin{gathered}
a'a''-bb=A,\quad aa''-b'b'=A',\quad aa'-b''b''=A''\\
b'b''-ab=B,\quad bb''-a'b'=B',\quad bb'-a''b''=B''\\
D=aa'a''+2bb'b''-abb-a'b'b'-a''b''b''
\end{gathered}
\]
so ist
\[
\begin{array}{rclll}
T&=&At+B''t'+B't''&\text{senkrecht zu}&t'\text{ und }t''\\
T'&=&B''t+A't'+Bt''&&t\text{ und }t''\\
T''&=&B't+Bt'+A''t''&&t\text{ und }t'
\end{array}
\]
und allgemein, wenn
\[
\lambda x+\lambda'x'+\lambda''x''=L
\]
die Gleichung einer Ebene ist, so wird die Linie
\[
\lambda T+\lambda'T'+\lambda''T''
\]
zu ihr senkrecht sein.

Ferner gilt
\[
\begin{aligned}
aT+b''T'+b'T''&=Dt\\
b''T+a'T'+bT''&=Dt'\\
b'T+bT'+a''T''&=Dt''
\end{aligned}
\]
und die Linien $t,t',t''$ sind senkrecht zu den Ebenen, deren Gleichungen
\[
\begin{aligned}
ax+b''x'+b'x''&=\mathrm{Const}\\
b''x+a'x'+bx''&=\mathrm{Const}\\
b'x+bx'+a''x''&=\mathrm{Const}
\end{aligned}
\]

\SourceLeaf{317}
Der doppelte Flächeninhalt des Dreiecks durch die Punkte $mt,m't',m''t''$ ist gleich der Quadratwurzel aus dem Wert der Form
\[
F\ldots\left(\begin{smallmatrix}A,&A',&A''\\B,&B',&B''\end{smallmatrix}\right)
\]
wenn man einsetzt: $X=m'm'',X'=mm'',X''=mm'$, während das sechsfache Volumen der Pyramide, die sich dadurch mit dem Punkt $0$ bildet, $=mm'm''\surd D$ wird, also ist das Lot
\[
=\surd\frac{D}{F(\frac1m,\frac1{m'},\frac1{m''})}
\]
$T,T',T''$ beziehen sich ebenso auf die Form $\left(\begin{smallmatrix}AD,&A'D,&A''D\\BD,&B'D,&B''D\end{smallmatrix}\right)$ wie $t,t',t''$ auf $\left(\begin{smallmatrix}a,&a',&a''\\b,&b',&b''\end{smallmatrix}\right)$

Die drei Wurzeln der Gleichung
\[
0=p^3-pp(a+a'+a'')+p(A+A'+A'')-D
\]
geben die Quadrate der drei Hauptachsen eines Ellipsoids an, das in dasjenige Parallelepiped eingeschrieben ist, auf welches sich die positive ternäre Form
\[
\left(\begin{smallmatrix}a,&a',&a''\\b,&b',&b''\end{smallmatrix}\right)
\quad\text{mit adjungierter Form }\left(\begin{smallmatrix}A,&A',&A''\\B,&B',&B''\end{smallmatrix}\right)
\quad\text{und Determinante }=-D
\]
bezieht.
\BookRule{.30\textwidth}

\subsection*{Beziehung der Raumverhältnisse auf ein gegebenes Tetraeder.}
\pdfbookmark[1]{Beziehung der Raumverhältnisse auf ein gegebenes Tetraeder}{ternary-tetraeder}
Es bezeichnen $(0),(1),(2),(3)$ die vier Ecken, gegenüberliegenden Flächen und Lote. Jeder Punkt des Raums $P$ hat dann bezüglich eines beliebigen Anfangspunkts $M$ vier Koordinaten $x,x',x'',x'''$, zwischen denen jedoch die Beziehung
\[
x+x'+x''+x'''=0
\]
gilt. Dabei bezeichnet $x$ den Quotienten, den man erhält, wenn man den Abstand des
\SourceLeaf{318}
Punktes $P$ von einer durch $M$ zur Ebene $(0)$ gelegten parallelen Ebene durch das Lot $(0)$ dividiert usw.

Allgemein ist dann
\[
\begin{aligned}
-(PM)^2={}&xx'(01)^2+xx''(02)^2+xx'''(03)^2\\
&+x'x''(12)^2+x'x'''(13)^2+x''x'''(23)^2
\end{aligned}
\]
\BookRule{.30\textwidth}

\emph{Das Grundgesetz der Kristallisation} lässt sich am kürzesten so aussprechen: Zwischen je fünf dabei vorkommenden Ebenen besteht folgende Beziehung:

Sind ihre Normalen auf der Kugeloberfläche $(0),(1),(2),(3),(4)$, so sind stets die Produkte $\sin102\pdot\sin304$, $\sin103\pdot\sin204$, $\sin203\pdot\sin104$ in einem rationalen Verhältnis; ist dies wie $\alpha:\beta:\gamma$, so ist $\beta=\alpha+\gamma$.

Sind die Koordinaten der $5$ Punkte auf der Kugeloberfläche
\[
\left.\begin{array}{ccc}
a&b&c\\a'&b'&c'\\a''&b''&c''\\a'''&b'''&c'''\\0&0&1
\end{array}\right\}
\quad\text{so müssen}\quad
\begin{gathered}
(ab'-ba')\pdot(a''b'''-b''a''')\\
(ab''-ba'')\pdot(a'''b'-b'''a')\\
(ab'''-ba''')\pdot(a'b''-b'a'')
\end{gathered}
\]
in rationalem Verhältnis stehen.

Allgemein seien $1,2,3,4,5$ die $5$ Punkte auf der Kugeloberfläche, $0$ der Mittelpunkt; dann stehen, wenn $12$ das Volumen des Tetraeders $0345$ bedeutet
\[
\left.\begin{array}{c}23\pdot45\\24\pdot53\\25\pdot34\end{array}\right\}
\quad\text{in rationalem Verhältnis}
\]
ebenso
\[
\left.\begin{array}{c}12\pdot34\\13\pdot42\\14\pdot23\end{array}\right\}\quad\text{usw.}
\]

